\section{Ход работы}
\begin{enumerate}
	\item select "Н\_ЛЮДИ"."ИМЯ", "Н\_СЕССИЯ"."УЧГОД" FROM "Н\_ЛЮДИ" RIGHT JOIN "Н\_СЕССИЯ" ON
	"ЧЛВК\_ИД" = "Н\_ЛЮДИ"."ИД" WHERE "Н\_ЛЮДИ"."ФАМИЛИЯ" > 'Ёлкин' AND "Н\_СЕССИЯ"."ИД" < 14 A
	ND "Н\_СЕССИЯ"."ИД" = 32199;
	\item select "Н\_ЛЮДИ"."ИМЯ", "Н\_ВЕДОМОСТИ"."ЧЛВК\_ИД", "Н\_СЕССИЯ"."УЧГОД" from "Н\_ЛЮДИ" left join "Н\_ВЕДОМОСТИ" on "Н\_ЛЮДИ"."ИД" = "Н\_ВЕДОМОСТИ"."ЧЛВК\_ИД" left join "Н\_СЕССИЯ"
	on "Н\_ЛЮДИ"."ИД" = "Н\_СЕССИЯ"."ЧЛВК\_ИД" where "Н\_ЛЮДИ"."ИД" > 152862 and "Н\_ВЕДОМОСТИ".
	"ДАТА" < '2010-06-18';
	
	\item select COUNT("Н\_ЛЮДИ"."ИД") as count from "Н\_ЛЮДИ" inner join "Н\_ОБУЧЕНИЯ" on "Н\_ЛЮДИ"."ИД" = "Н\_ОБУЧЕНИЯ"."ЧЛВК\_ИД" inner join "Н\_УЧЕНИКИ" on "Н\_УЧЕНИКИ"."ЧЛВК\_ИД" = "Н\_ОБУЧЕНИЯ"."ЧЛВК\_ИД" inner join "Н\_ПЛАНЫ" on "Н\_УЧЕНИКИ"."ПЛАН\_ИД" = "Н\_ПЛАНЫ"."ИД" inner join "Н\_ФОРМЫ\_ОБУЧЕНИЯ" on "Н\_ФОРМЫ\_ОБУЧЕНИЯ"."ИД" = "Н\_ПЛАНЫ"."ФО\_ИД" where "Н\_ЛЮДИ"."ДАТА\_РОЖДЕНИЯ" > '2026-03-15'::DATE - INTERVAL '20 years';
	
	\item SELECT "Н\_ГРУППЫ\_ПЛАНОВ"."ПЛАН\_ИД", COUNT(distinct "Н\_ГРУППЫ\_ПЛАНОВ"."ГРУППА") as cnt from "Н\_ГРУППЫ\_ПЛАНОВ" inner join "Н\_ПЛАНЫ" on "Н\_ПЛАНЫ"."ПЛАН\_ИД" = "Н\_ГРУППЫ\_ПЛАНОВ"."ПЛАН\_ИД" inner join "Н\_ОТДЕЛЫ" on "Н\_ОТДЕЛЫ"."ИД" = "Н\_ПЛАНЫ"."ОТД\_ИД" where "Н\_ОТДЕЛЫ"."КОРОТКОЕ\_ИМЯ" = 'ВТ' group by "Н\_ГРУППЫ\_ПЛАНОВ"."ПЛАН\_ИД" having COUNT("Н\_ГРУППЫ\_ПЛАНОВ"."ГРУППА") < 2;
	
	\item \begin{lstlisting}[language=SQL] 
		SELECT
	"Н\_УЧЕНИКИ"."ИД",
	"Н\_ЛЮДИ"."ФАМИЛИЯ",
	"Н\_ЛЮДИ"."ИМЯ",
	"Н\_ЛЮДИ"."ОТЧЕСТВО",
	AVG(CAST("Н\_ВЕДОМОСТИ"."ОЦЕНКА" AS INTEGER)) as avg\_grade
	FROM "Н\_УЧЕНИКИ"
	INNER JOIN "Н\_ОБУЧЕНИЯ" ON "Н\_ОБУЧЕНИЯ"."ЧЛВК\_ИД" = "Н\_УЧЕНИКИ"."ЧЛВК\_ИД"
	INNER JOIN "Н\_ЛЮДИ" ON "Н\_ЛЮДИ"."ИД" = "Н\_ОБУЧЕНИЯ"."ЧЛВК\_ИД"
	INNER JOIN "Н\_ВЕДОМОСТИ" ON "Н\_ВЕДОМОСТИ"."ЧЛВК\_ИД" = "Н\_ЛЮДИ"."ИД" 
	WHERE
	"Н\_УЧЕНИКИ"."ГРУППА" = '4100'
	AND "Н\_ВЕДОМОСТИ"."ОЦЕНКА" ~ '\^[0-9]+\$'
	GROUP BY
	"Н\_УЧЕНИКИ"."ИД",
	"Н\_ЛЮДИ"."ФАМИЛИЯ",
	"Н\_ЛЮДИ"."ИМЯ",
	"Н\_ЛЮДИ"."ОТЧЕСТВО"
	HAVING
	AVG(CAST("Н\_ВЕДОМОСТИ"."ОЦЕНКА" AS INTEGER)) = (
	SELECT
	MIN(avg\_grade)
	FROM (
	SELECT
	AVG(CAST("Н\_ВЕДОМОСТИ"."ОЦЕНКА" AS INTEGER)) as avg\_grade
	FROM "Н\_УЧЕНИКИ"
	INNER JOIN "Н\_ОБУЧЕНИЯ" ON "Н\_ОБУЧЕНИЯ"."ЧЛВК\_ИД" = "Н\_УЧЕНИКИ"."ЧЛВК\_ИД"
	INNER JOIN "Н\_ЛЮДИ" ON "Н\_ЛЮДИ"."ИД" = "Н\_ОБУЧЕНИЯ"."ЧЛВК\_ИД"
	INNER JOIN "Н\_ВЕДОМОСТИ" ON "Н\_ВЕДОМОСТИ"."ЧЛВК\_ИД" = "Н\_ЛЮДИ"."ИД"
	WHERE
	"Н\_УЧЕНИКИ"."ГРУППА" = '1100'
	AND "Н\_ВЕДОМОСТИ"."ОЦЕНКА" ~ '\^[0-9]+\$'
	GROUP BY "Н\_УЧЕНИКИ"."ЧЛВК\_ИД"
	));
	\end{lstlisting}
	
	
	\item select "Н\_УЧЕНИКИ"."ГРУППА",
	"Н\_УЧЕНИКИ"."ЧЛВК\_ИД",
	"Н\_ЛЮДИ"."ФАМИЛИЯ",
	"Н\_ЛЮДИ"."ИМЯ",
	"Н\_ЛЮДИ"."ОТЧЕСТВО",
	"Н\_УЧЕНИКИ"."П\_ПРКОК\_ИД",
	"Н\_УЧЕНИКИ"."СОСТОЯНИЕ"
	from "Н\_УЧЕНИКИ"
	inner join "Н\_ЛЮДИ"
	on "Н\_ЛЮДИ"."ИД" = "Н\_УЧЕНИКИ"."ЧЛВК\_ИД"
	inner join "Н\_ПЛАНЫ"
	on "Н\_УЧЕНИКИ"."ПЛАН\_ИД" = "Н\_ПЛАНЫ"."ИД"
	inner join "Н\_ФОРМЫ\_ОБУЧЕНИЯ"
	on "Н\_ФОРМЫ\_ОБУЧЕНИЯ"."ИД" = "Н\_ПЛАНЫ"."ФО\_ИД"
	where "Н\_УЧЕНИКИ"."НАЧАЛО" > '2012-09-01'
	and "Н\_ФОРМЫ\_ОБУЧЕНИЯ"."НАИМЕНОВАНИЕ" = 'Заочная'
	and "Н\_ПЛАНЫ"."КУРС" = 1
	and exists (
	select 1
	from "Н\_НАПРАВЛЕНИЯ\_СПЕЦИАЛ"
	inner join "Н\_НАПР\_СПЕЦ"
	on "Н\_НАПР\_СПЕЦ"."ИД" = "Н\_НАПРАВЛЕНИЯ\_СПЕЦИАЛ"."НАПС\_ИД"
	where "Н\_НАПРАВЛЕНИЯ\_СПЕЦИАЛ"."ИД" = "Н\_ПЛАНЫ"."НАПС\_ИД"
	and "Н\_НАПР\_СПЕЦ"."НАИМЕНОВАНИЕ" = 'Программная инженерия'
	);
	
	
	
	\item select count(*) from (SELECT
	"Н\_УЧЕНИКИ"."ЧЛВК\_ИД",
	AVG(CAST("Н\_ВЕДОМОСТИ"."ОЦЕНКА" AS INTEGER)) as avg\_grade
	FROM "Н\_УЧЕНИКИ"
	INNER JOIN "Н\_ВЕДОМОСТИ" ON "Н\_УЧЕНИКИ"."ЧЛВК\_ИД" = "Н\_ВЕДОМОСТИ"."Ч
	ЛВК\_ИД"
	WHERE "Н\_ВЕДОМОСТИ"."ОЦЕНКА" ~ '\^[0-9]+\$'
	GROUP BY "Н\_УЧЕНИКИ"."ЧЛВК\_ИД"
	HAVING
	MAX(CAST("Н\_ВЕДОМОСТИ"."ОЦЕНКА" AS INTEGER)) = 5
	AND MIN(CAST("Н\_ВЕДОМОСТИ"."ОЦЕНКА" AS INTEGER)) = 5);
\end{enumerate}

Вывод: углубил свои знания в SQL командах. Понял, что LaTex  так себе текстовый редактор для выполнения отчета по БД.